\documentclass[12pt,a4paper]{article}

% ─── Paquetes ────────────────────────────────────────────────────────────────
\usepackage[utf8]{inputenc}
\usepackage[T1]{fontenc}
\usepackage[spanish,mexico]{babel}
\usepackage[margin=2.5cm]{geometry}
\usepackage{graphicx}
\usepackage{xcolor}
\usepackage{hyperref}
\usepackage{booktabs}
\usepackage{longtable}
\usepackage{array}
\usepackage{enumitem}
\usepackage{fancyhdr}
\usepackage{titlesec}
\usepackage{listings}
\usepackage{tcolorbox}
\usepackage{tabularx}
\usepackage{multirow}
\usepackage{amsmath}
\usepackage{lmodern}
\usepackage{microtype}
\setlength{\parskip}{6pt}
\setlength{\parindent}{0pt}

% ─── Colores corporativos (tomados de styles.css) ────────────────────────────
\definecolor{ink900}{HTML}{0b1736}
\definecolor{sky500}{HTML}{1c78c0}
\definecolor{teal500}{HTML}{0f8d8f}
\definecolor{lime500}{HTML}{5f9d2f}
\definecolor{amber500}{HTML}{c97a1a}
\definecolor{rose500}{HTML}{c94a56}
\definecolor{paper100}{HTML}{f8fafc}
\definecolor{graylight}{HTML}{e2e8f0}

% ─── Configuración de hipervínculos ─────────────────────────────────────────
\hypersetup{
  colorlinks=true,
  linkcolor=ink900,
  urlcolor=sky500,
  citecolor=teal500,
  pdftitle={Servicia -- Documentacion Tecnica y Funcional},
  pdfauthor={Direccion General -- Batia},
  pdfsubject={Plataforma de Aseguramiento Contable}
}

% ─── Encabezado y pie de página ─────────────────────────────────────────────
\pagestyle{fancy}
\fancyhf{}
\fancyhead[L]{\small\color{ink900}\textbf{Servicia} $\cdot$ Aseguramiento Contable}
\fancyhead[R]{\small\color{sky500}Batia $\cdot$ Piloto}
\fancyfoot[C]{\small\thepage}
\fancyfoot[R]{\small\color{gray}Confidencial}
\renewcommand{\headrulewidth}{0.4pt}
\renewcommand{\footrulewidth}{0.2pt}

% ─── Formato de títulos ──────────────────────────────────────────────────────
\titleformat{\section}
  {\large\bfseries\color{ink900}}
  {\thesection.}{0.6em}{}[\color{sky500}\titlerule]

\titleformat{\subsection}
  {\normalsize\bfseries\color{sky500}}
  {\thesubsection.}{0.5em}{}

\titleformat{\subsubsection}
  {\small\bfseries\color{teal500}}
  {\thesubsubsection.}{0.4em}{}

% ─── Cajas de información ────────────────────────────────────────────────────
\tcbuselibrary{skins, breakable}

\newtcolorbox{infobox}[1][]{
  enhanced, breakable,
  colback=paper100, colframe=sky500,
  fonttitle=\bfseries\color{white},
  colbacktitle=sky500,
  title=#1, left=4pt, right=4pt
}

\newtcolorbox{alertbox}[1][]{
  enhanced, breakable,
  colback=amber500!10, colframe=amber500,
  fonttitle=\bfseries\color{white},
  colbacktitle=amber500,
  title=#1, left=4pt, right=4pt
}

\newtcolorbox{rulebox}[1][]{
  enhanced, breakable,
  colback=rose500!8, colframe=rose500,
  fonttitle=\bfseries\color{white},
  colbacktitle=rose500,
  title=#1, left=4pt, right=4pt
}

% ─── Configuración de listings ───────────────────────────────────────────────
\lstset{
  basicstyle=\ttfamily\footnotesize,
  keywordstyle=\color{sky500}\bfseries,
  commentstyle=\color{teal500}\itshape,
  stringstyle=\color{lime500},
  numberstyle=\tiny\color{gray},
  numbers=left, numbersep=6pt,
  frame=single, framesep=4pt,
  rulecolor=\color{ink900!20},
  backgroundcolor=\color{paper100},
  breaklines=true, breakatwhitespace=true,
  showspaces=false, showstringspaces=false,
  tabsize=4, captionpos=b,
  xleftmargin=10pt
}

\lstdefinelanguage{Python}{
  keywords={def,class,return,import,from,if,else,elif,for,while,with,as,
            try,except,finally,pass,True,False,None,and,or,not,in,is,
            lambda,yield,global,nonlocal,del,raise,assert,bool},
  sensitive=true,
  comment=[l]{\#},
  morestring=[b]",
  morestring=[b]'
}

\lstdefinestyle{plain}{
  language={},
  numbers=none,
  frame=single,
  backgroundcolor=\color{ink900!5},
  basicstyle=\ttfamily\footnotesize
}

% ─── Documento ───────────────────────────────────────────────────────────────
\begin{document}

% ═══════════════════════════════════════════════════════════════════════════════
%  PORTADA
% ═══════════════════════════════════════════════════════════════════════════════
\begin{titlepage}
  \pagecolor{ink900}
  \color{white}
  \centering
  \vspace*{3cm}

  {\fontsize{48}{52}\selectfont\bfseries SERVICIA}\\[0.4cm]
  {\large\color{sky500!80!white} Plataforma de Aseguramiento Contable}\\[0.2cm]
  {\normalsize\color{graylight} Piloto Batia}

  \vspace{2cm}
  \noindent\color{sky500}\rule{\linewidth}{1pt}
  \vspace{0.5cm}

  {\large Documentaci\'on T\'ecnica y Funcional}\\[0.3cm]
  {\normalsize An\'alisis de arquitectura, flujos operativos, modelos de datos\\
  e instrucciones para Unidad de Inteligencia y Contabilidad}

  \vfill

  \begin{tabular}{ll}
    \textbf{Emitido por:}    & Direcci\'on General \\[4pt]
    \textbf{Clasificaci\'on:}  & Confidencial $\cdot$ Uso interno \\[4pt]
    \textbf{Versi\'on:}        & 2.0 \\[4pt]
    \textbf{Fecha:}          & Abril 2026 \\[4pt]
    \textbf{Stack:}          & Flask 3.0 $\cdot$ SQLAlchemy $\cdot$ PostgreSQL $\cdot$ Cloud Run \\
  \end{tabular}

  \vspace{1cm}
  \noindent\color{sky500!50!white}\rule{\linewidth}{0.4pt}
  \vspace{0.3cm}

  {\small\color{graylight}
    Destinatarios: Unidad de Inteligencia (UI) $\cdot$ \'Area de Contabilidad $\cdot$ Tesorer\'ia $\cdot$ Direcci\'on}
\end{titlepage}

\pagecolor{white}
\color{black}

% ═══════════════════════════════════════════════════════════════════════════════
%  TABLA DE CONTENIDOS
% ═══════════════════════════════════════════════════════════════════════════════
\tableofcontents
\newpage

% ═══════════════════════════════════════════════════════════════════════════════
\section{Resumen Ejecutivo}
% ═══════════════════════════════════════════════════════════════════════════════

\begin{infobox}[Descripci\'on General]
\textbf{Servicia} es una plataforma web de aseguramiento contable desarrollada en Flask
(Python) para \textbf{Batia}, empresa de gesti\'on de servicios integrales. El sistema
digitaliza y controla el ciclo completo de vida de cada contrataci\'on registrada en SINGA:
desde la apertura del expediente hasta el cierre fiscal verificado por Contabilidad,
pasando por validaciones SAT/EFOS, traspasos SPEI, vinculaci\'on de CFDI y conciliaci\'on
bancaria.
\end{infobox}

\bigskip

El proyecto responde a una necesidad cr\'itica: garantizar que \textbf{toda operaci\'on de
pago est\'e completamente materializada}, documentada y conciliada, reduciendo el riesgo
fiscal ante el SAT y permitiendo que cualquier auditor\'ia interna o externa pueda verificar
en minutos la legalidad y trazabilidad de cada transacci\'on.

\subsection*{Caracter\'isticas principales}

\begin{itemize}[leftmargin=1.5em]
  \item \textbf{Flujo operativo de 8 pasos} con bloqueos secuenciales: ning\'un paso puede
    ejecutarse sin que el anterior est\'e completado.
  \item \textbf{Gesti\'on de folios SINGA} con checklist de entregables diferenciado seg\'un
    tipo de proveedor (outsourcing, materiales, servicios).
  \item \textbf{Control SAT/EFOS}: bloquea pagos a proveedores sin Opini\'on de
    Cumplimiento vigente.
  \item \textbf{M\'odulo de conciliaci\'on bancaria} con carga de PDF y extracci\'on
    autom\'atica de movimientos.
  \item \textbf{Sistema de alertas autom\'aticas} (dep\'osito sin folio, monto excedido,
    cuenta no autorizada, diferencia CFDI-banco).
  \item \textbf{Tablero ejecutivo} personalizado por rol con KPIs en tiempo real.
  \item \textbf{Despliegue en Google Cloud Run} con base de datos PostgreSQL gestionada.
\end{itemize}

\newpage

% ═══════════════════════════════════════════════════════════════════════════════
\section{Arquitectura del Sistema}
% ═══════════════════════════════════════════════════════════════════════════════

\subsection{Stack Tecnol\'ogico}

\begin{table}[h!]
\centering
\begin{tabularx}{\textwidth}{>{\bfseries}lX}
\toprule
Capa & Tecnolog\'ia \\
\midrule
Lenguaje & Python 3.11 \\
Framework Web & Flask 3.0.3 \\
ORM & Flask-SQLAlchemy 3.1.1 \\
Migraciones & Flask-Migrate 4.0.7 (Alembic) \\
Autenticaci\'on & Flask-JWT-Extended 4.6.0 + sesiones de servidor \\
Templates & Jinja2 (server-side rendering) \\
BD Desarrollo & SQLite (\texttt{instance/servicia.db}) \\
BD Producci\'on & PostgreSQL (Cloud SQL, GCP) \\
Servidor WSGI & Gunicorn 23.0.0 (1 worker, 8 threads, timeout 120\,s) \\
Contenedor & Docker (Python 3.11-slim) \\
Plataforma Cloud & Google Cloud Run + Cloud Build \\
Excel & openpyxl 3.1.5 \\
PDF parsing & pypdf 5.4.0 \\
\bottomrule
\end{tabularx}
\caption{Stack tecnol\'ogico de Servicia}
\end{table}

\subsection{Estructura del Proyecto}

El proyecto sigue el patr\'on \textit{Application Factory} de Flask con
\textit{blueprints} modulares. Cada \'area funcional tiene su propio blueprint
con sus rutas y plantillas.

\begin{lstlisting}[style=plain, caption={\'Arbol de directorios del proyecto}]
contabilidad/
  app/
    __init__.py          <- Application factory
    config.py            <- Variables de entorno
    extensions.py        <- Registro de db, migrate, jwt
    models.py            <- 18 modelos SQLAlchemy
    audit.py             <- Log de auditoria
    seed.py              <- Datos demo idempotentes
    auth/routes.py       <- Login, logout, JWT
    main/routes.py       <- Landing, dashboard, busqueda
    folios/routes.py     <- CRUD de folios SINGA
    operations/
      routes.py          <- SAT, CFDI, traspasos, conciliacion
      services.py        <- Logica de negocio y workflow
    static/
      css/styles.css     <- Hoja de estilos
      js/main.js         <- Interacciones frontend
    templates/           <- Plantillas Jinja2 por modulo
  instance/servicia.db   <- BD SQLite (desarrollo)
  ref/                   <- Manual de usuario y referencias
  uploads/               <- Archivos subidos por usuarios
  run.py                 <- Entrypoint desarrollo
  wsgi.py                <- Entrypoint produccion
  Dockerfile
  cloudbuild.yaml        <- Pipeline CI/CD en GCP
  requirements.txt
\end{lstlisting}

\subsection{Application Factory y Blueprints}

La funci\'on \texttt{create\_app()} en \texttt{app/\_\_init\_\_.py} sigue el patr\'on
\textit{factory} para permitir m\'ultiples instancias (testing, desarrollo, producci\'on):

\begin{lstlisting}[language=Python, caption={Registro de blueprints en create\_app()}]
def create_app(config_class=Config):
    app = Flask(__name__)
    app.config.from_object(config_class)

    # Inicializacion de extensiones
    db.init_app(app)
    migrate.init_app(app, db)
    jwt.init_app(app)

    # Registro de blueprints
    app.register_blueprint(auth_bp,       url_prefix='/auth')
    app.register_blueprint(main_bp)
    app.register_blueprint(folios_bp,     url_prefix='/folios')
    app.register_blueprint(operations_bp, url_prefix='/operacion')

    # Headers de seguridad en cada respuesta
    @app.after_request
    def security_headers(response):
        response.headers['X-Content-Type-Options'] = 'nosniff'
        response.headers['X-Frame-Options'] = 'DENY'
        return response

    return app
\end{lstlisting}

\subsection{Seguridad y Autenticaci\'on}

El sistema implementa autenticaci\'on dual:

\begin{itemize}
  \item \textbf{Sesiones de servidor} (Flask session): para el acceso v\'ia navegador.
    El decorador \texttt{@login\_required} verifica la presencia de \texttt{user\_id}
    en sesi\'on.
  \item \textbf{JWT (JSON Web Tokens)}: endpoint de API
    \texttt{POST /auth/api/token}, que emite tokens con claims de rol y nombre.
\end{itemize}

\begin{lstlisting}[language=Python, caption={Decorador de control de acceso por rol}]
def roles_required(*allowed_roles):
    def decorator(f):
        @wraps(f)
        def wrapper(*args, **kwargs):
            role = session.get('role')
            if role not in allowed_roles:
                flash('No tienes permisos para esta accion.', 'error')
                return redirect(url_for('main.dashboard'))
            return f(*args, **kwargs)
        return wrapper
    return decorator
\end{lstlisting}

\newpage

% ═══════════════════════════════════════════════════════════════════════════════
\section{Modelo de Datos}
% ═══════════════════════════════════════════════════════════════════════════════

La base de datos contiene \textbf{18 modelos SQLAlchemy} que se mapean a las tablas
de la BD. A continuaci\'on se describen los m\'as importantes.

\subsection{Entidades Principales}

\subsubsection{Role y User}

\begin{lstlisting}[language=Python, caption={Modelo User con hash de contrase\~na}]
class User(db.Model):
    __tablename__ = 'users'
    id            = db.Column(db.Integer, primary_key=True)
    full_name     = db.Column(db.String(120), nullable=False)
    email         = db.Column(db.String(120), unique=True, nullable=False)
    password_hash = db.Column(db.String(256), nullable=False)
    role_id       = db.Column(db.Integer, db.ForeignKey('roles.id'))
    is_active     = db.Column(db.Boolean, default=True)
    role          = db.relationship('Role', lazy='joined')

    def set_password(self, raw):
        self.password_hash = generate_password_hash(raw)

    def check_password(self, raw):
        return check_password_hash(self.password_hash, raw)
\end{lstlisting}

Los roles disponibles son: \texttt{direccion}, \texttt{ui}, \texttt{contabilidad},
\texttt{tesoreria} y \texttt{auditor}.

\subsubsection{Provider y AuthorizedAccount}

\begin{lstlisting}[language=Python, caption={Modelos Provider y AuthorizedAccount}]
class Provider(db.Model):
    __tablename__ = 'providers'
    id            = db.Column(db.Integer, primary_key=True)
    name          = db.Column(db.String(200), unique=True, nullable=False)
    provider_type = db.Column(db.String(30))  # outsourcing, materiales, servicios
    sat_status    = db.Column(db.String(30))
    sat_valid_until = db.Column(db.Date)

class AuthorizedAccount(db.Model):
    __tablename__ = 'authorized_accounts'
    clabe     = db.Column(db.String(18), unique=True, nullable=False)
    rfc       = db.Column(db.String(13))
    is_active = db.Column(db.Boolean, default=True)
    __table_args__ = (
        db.CheckConstraint("LENGTH(clabe) = 18"),)
\end{lstlisting}

\subsubsection{Folio --- El Expediente Central}

El modelo \texttt{Folio} representa cada contrataci\'on registrada en SINGA. Es la
entidad pivote: todos los dem\'as registros operativos apuntan a un folio.

\begin{lstlisting}[language=Python, caption={Modelo Folio con estatus sem\'aforo}]
class Folio(db.Model):
    __tablename__ = 'folios'
    id                        = db.Column(db.Integer, primary_key=True)
    singa_number              = db.Column(db.String(50), unique=True, nullable=False)
    provider_id               = db.Column(db.Integer, db.ForeignKey('providers.id'))
    provider_type             = db.Column(db.String(30))
    communication_responsible = db.Column(db.String(120))
    contract_amount           = db.Column(db.Numeric(14, 2))
    budget_amount             = db.Column(db.Numeric(14, 2))
    status                    = db.Column(db.String(30), default='pendiente')
    created_at                = db.Column(db.DateTime, default=datetime.utcnow)

    # status: pendiente | en_proceso | listo_para_cierre | cerrado | critico
    __table_args__ = (
        db.CheckConstraint("status IN ('pendiente','en_proceso',"
            "'listo_para_cierre','cerrado','alerta','critico')"),
        db.CheckConstraint("budget_amount > 0"),
    )
\end{lstlisting}

\subsubsection{FolioWorkflow --- Los 8 Pasos}

\begin{lstlisting}[language=Python, caption={Modelo FolioWorkflow}]
class FolioWorkflow(db.Model):
    __tablename__ = 'folio_workflows'
    folio_id                 = db.Column(db.Integer,
                                 db.ForeignKey('folios.id'), unique=True)
    step_1_ui_folio          = db.Column(db.Boolean, default=False)
    step_2_contab_sat        = db.Column(db.Boolean, default=False)
    step_3_ui_deliverables   = db.Column(db.Boolean, default=False)
    step_4_tesoreria_transfer= db.Column(db.Boolean, default=False)
    step_5_contab_cfdi       = db.Column(db.Boolean, default=False)
    step_6_contab_reconciliation = db.Column(db.Boolean, default=False)
    step_7_ui_ready_to_close = db.Column(db.Boolean, default=False)
    step_8_contab_closed     = db.Column(db.Boolean, default=False)
    updated_at               = db.Column(db.DateTime, onupdate=datetime.utcnow)
\end{lstlisting}

\subsection{Entidades de Control Fiscal}

\subsubsection{ProviderCompliance}

Almacena el estado de cumplimiento SAT/EFOS de cada proveedor con relaci\'on 1:1.

\begin{table}[h!]
\centering
\begin{tabularx}{\textwidth}{>{\ttfamily\small}l>{\small}X}
\toprule
\textnormal{\textbf{Campo}} & \textbf{Descripci\'on} \\
\midrule
sat\_opinion\_status & \texttt{no\_cargada} / \texttt{vigente} / \texttt{vencida} \\
sat\_valid\_until    & Fecha de vencimiento de la opini\'on SAT \\
efos\_flag           & Aparece en lista EFOS del SAT (booleano) \\
activity\_match      & Actividad econ\'omica congruente con el servicio \\
rfc\_matches\_account & RFC coincide con CLABE registrada \\
notes                & Observaciones de Contabilidad \\
\bottomrule
\end{tabularx}
\caption{Campos de ProviderCompliance}
\end{table}

\subsubsection{CfdiRecord}

\begin{lstlisting}[language=Python, caption={Modelo CfdiRecord}]
class CfdiRecord(db.Model):
    __tablename__ = 'cfdi_records'
    uuid      = db.Column(db.String(36), unique=True, nullable=False)
    xml_ref   = db.Column(db.String(255))   # ruta del XML subido
    pdf_ref   = db.Column(db.String(255))   # ruta del PDF subido
    amount    = db.Column(db.Numeric(14, 2))
    issued_at = db.Column(db.Date)
    status    = db.Column(db.String(20), default='vigente')
    # status: vigente | cancelado
\end{lstlisting}

\subsubsection{BankMovement y ReconciliationAlert}

\begin{table}[h!]
\centering
\begin{tabularx}{\textwidth}{>{\ttfamily\small}l>{\small}X}
\toprule
\multicolumn{2}{l}{\textbf{ReconciliationAlert --- Tipos de alerta autom\'atica}} \\
\midrule
deposito\_sin\_folio     & Dep\'osito bancario sin folio SINGA asociado \\
excede\_presupuesto      & Monto de traspaso o movimiento supera el presupuesto \\
cuenta\_no\_autorizada   & Movimiento a cuenta no registrada en la plataforma \\
diferencia\_cfdi\_banco  & Diferencia entre CFDI registrado y movimiento bancario \\
\midrule
\multicolumn{2}{l}{\textbf{Severidades:} \texttt{media} / \texttt{alta}} \\
\multicolumn{2}{l}{\textbf{Estatus:} \texttt{abierta} / \texttt{resuelta}} \\
\bottomrule
\end{tabularx}
\caption{Tipos de alerta de conciliaci\'on}
\end{table}

\subsection{Diagrama de Relaciones}

\begin{lstlisting}[style=plain, caption={Relaciones entre entidades principales}]
Role          1 ---< N   User
Provider      1 ---< N   AuthorizedAccount
Provider      1 ---< N   Folio
Provider      1 ---  1   ProviderCompliance
Folio         1 ---  1   FolioWorkflow
Folio         1 ---< N   FolioDeliverableItem
Folio         1 ---< N   Transfer
Folio         1 ---< N   CfdiRecord
Folio         1 ---< N   BankMovement
Folio         1 ---< N   ReconciliationAlert
Folio         1 ---< N   BudgetConfig
AuthorizedAccount  1 ---< N  Transfer
EvidenceFile  N ---  1   Folio
EvidenceFile  N ---  1   FolioDeliverableItem (opcional)
AuditLog      N ---  1   User (nullable)
MonthlyDirectionReport  N ---  1  User
\end{lstlisting}

\newpage

% ═══════════════════════════════════════════════════════════════════════════════
\section{Flujo Operativo de 8 Pasos}
% ═══════════════════════════════════════════════════════════════════════════════

\begin{infobox}[Principio fundamental]
\textbf{Ning\'un paso puede saltarse.} El sistema aplica bloqueos secuenciales en
c\'odigo: intentar registrar un traspaso sin que SAT est\'e validado genera un error
de negocio inmediato. El workflow avanza solo cuando el paso anterior est\'a completo.
\end{infobox}

\bigskip

\begin{longtable}{>{\bfseries}cp{2.5cm}p{6cm}p{3.5cm}}
\toprule
\# & \'Area & Acci\'on en la plataforma & Desbloquea \\
\midrule
\endfirsthead
\multicolumn{4}{c}{\small\textit{(continuaci\'on)}} \\
\toprule
\# & \'Area & Acci\'on & Desbloquea \\
\midrule
\endhead
\bottomrule
\endfoot
1 & UI & Registrar folio SINGA con proveedor, tipo y responsable & Apertura del expediente \\[6pt]
2 & Contabilidad & Verificar proveedor en SAT y subir Opini\'on de Cumplimiento & Cuenta autorizada para traspaso \\[6pt]
3 & UI & Subir entregables operativos seg\'un tipo de proveedor & Expediente pasa a ``En proceso'' \\[6pt]
4 & Tesorer\'ia & Registrar traspaso SPEI con folio bancario & Conciliaci\'on del monto \\[6pt]
5 & Contabilidad & Vincular CFDI al folio y validar monto vs presupuesto & Sistema verifica excedente \\[6pt]
6 & Contabilidad & Conciliar con estado de cuenta PDF & Sem\'aforo del Dashboard se actualiza \\[6pt]
7 & UI & Confirmar que todos los entregables est\'an completos & Folio pasa a ``Listo para cierre'' \\[6pt]
8 & Contabilidad & Revisar y aprobar expediente completo & Folio marcado ``Cerrado'' \\
\end{longtable}

\subsection{L\'ogica de Rec\'alculo Autom\'atico}

La funci\'on \texttt{recalculate\_workflow\_from\_data()} en \texttt{services.py}
implementa las reglas de avance de cada paso:

\begin{lstlisting}[language=Python, caption={Recalculo del workflow -- services.py}]
def recalculate_workflow_from_data(folio, workflow, deliverables,
                                   has_transfer, has_cfdi,
                                   has_reconciliation):
    # Paso 1: siempre True (folio ya existe)
    workflow.step_1_ui_folio = True

    # Paso 2: SAT/EFOS vigente
    compliance = ProviderCompliance.query.filter_by(
        provider_id=folio.provider_id).first()
    workflow.step_2_contab_sat = (
        compliance and
        compliance.sat_opinion_status == 'vigente' and
        not compliance.efos_flag
    )

    # Paso 3: todos los entregables UI subidos
    ui_items = [d for d in deliverables if d.owner_area == 'ui']
    workflow.step_3_ui_deliverables = (
        bool(ui_items) and all(d.uploaded for d in ui_items)
    )

    # Paso 4: traspaso solo si SAT validado
    workflow.step_4_tesoreria_transfer = (
        workflow.step_2_contab_sat and has_transfer
    )

    # Paso 5: CFDI solo si traspaso realizado
    workflow.step_5_contab_cfdi = (
        workflow.step_4_tesoreria_transfer and has_cfdi
    )

    # Paso 6: conciliacion solo si CFDI vigente
    workflow.step_6_contab_reconciliation = (
        workflow.step_5_contab_cfdi and has_reconciliation
    )

    # Paso 7: listo cierre si UI completo y conciliacion completa
    workflow.step_7_ui_ready_to_close = (
        workflow.step_3_ui_deliverables and
        workflow.step_6_contab_reconciliation
    )

    # Paso 8: cerrado si listo + conciliacion
    workflow.step_8_contab_closed = (
        workflow.step_7_ui_ready_to_close and
        workflow.step_6_contab_reconciliation
    )
\end{lstlisting}

\subsection{Estatus Sem\'aforo del Folio}

\begin{table}[h!]
\centering
\begin{tabularx}{\textwidth}{>{\ttfamily}lX}
\toprule
\textbf{Estatus} & \textbf{Condici\'on} \\
\midrule
critico             & Existen alertas de severidad ``alta'' abiertas \\
cerrado             & Paso 8 completo \\
listo\_para\_cierre & Pasos 7 y 6 completos \\
en\_proceso         & Alguno de los pasos 3, 4 o 5 completo \\
pendiente           & Estado inicial (ning\'un paso avanzado) \\
\bottomrule
\end{tabularx}
\caption{Estatus sem\'aforo de los folios SINGA}
\end{table}

\newpage

% ═══════════════════════════════════════════════════════════════════════════════
\section{M\'odulos del Sistema}
% ═══════════════════════════════════════════════════════════════════════════════

\subsection{M\'odulo de Autenticaci\'on (\texttt{/auth})}

\begin{itemize}
  \item \textbf{GET/POST \texttt{/auth/login}}: Valida email y contrase\~na (bcrypt v\'ia
    Werkzeug), crea sesi\'on con \texttt{user\_id}, \texttt{user\_name} y \texttt{role}.
    Registra evento en \texttt{AuditLog}.
  \item \textbf{POST \texttt{/auth/logout}}: Limpia sesi\'on y redirige al landing.
  \item \textbf{POST \texttt{/auth/api/token}}: Emite JWT con claims \texttt{role} y
    \texttt{name} para consumo de API externa.
\end{itemize}

\subsection{M\'odulo Principal (\texttt{/})}

\subsubsection{Dashboard Ejecutivo}

El dashboard se personaliza seg\'un el rol del usuario autenticado:

\begin{table}[h!]
\centering
\begin{tabularx}{\textwidth}{>{\bfseries}lX}
\toprule
Rol & Informaci\'on en el Dashboard \\
\midrule
UI & Entregables pendientes, folios activos, folios cr\'iticos \\
Contabilidad & Entregables contables, alertas abiertas, folios sin CFDI \\
Tesorer\'ia & Folios por traspasar, alertas de severidad alta \\
Direcci\'on & Reportes mensuales, alertas antiguas (+5 d\'ias), acuses firmados \\
Auditor & Vista completa de KPIs sin permisos de escritura \\
\bottomrule
\end{tabularx}
\caption{Informaci\'on personalizada por rol en el dashboard}
\end{table}

\subsubsection{B\'usqueda Global (\texttt{/busqueda})}

Permite buscar por texto libre. Devuelve hasta 25 resultados por categor\'ia:
folios (SINGA, responsable, proveedor), proveedores, UUIDs de CFDI, referencias
SPEI y alertas.

\subsection{M\'odulo de Folios (\texttt{/folios})}

Vista \textit{master-detail}: lista de folios con filtros en el panel izquierdo
y detalle del folio seleccionado en el panel derecho (workflow, entregables,
transferencias, CFDIs, movimientos, alertas, evidencias y eventos).

\subsubsection{Creaci\'on de un Folio}

Al crear un nuevo folio, el sistema:
\begin{enumerate}
  \item Valida que el n\'umero SINGA sea \'unico.
  \item Verifica que el proveedor tenga al menos una \textbf{cuenta bancaria autorizada
    activa}. Si no la tiene, bloquea la operaci\'on.
  \item Valida que el presupuesto no supere el monto del contrato.
  \item Crea el registro \texttt{FolioWorkflow} (8 pasos en falso).
  \item Genera autom\'aticamente el checklist de \texttt{FolioDeliverableItem}
    seg\'un el tipo de proveedor.
\end{enumerate}

\subsection{M\'odulo Operativo (\texttt{/operacion})}

\subsubsection{Cumplimiento SAT/EFOS}

Contabilidad gestiona el estado fiscal de cada proveedor. Al actualizar:
\begin{itemize}
  \item Se registra la Opini\'on de Cumplimiento con su fecha de vencimiento.
  \item Se marcan flags: EFOS, congruencia de actividad, coincidencia RFC-CLABE.
  \item El sistema \textbf{recalcula autom\'aticamente el estatus de todos los folios}
    activos del proveedor.
\end{itemize}

\begin{rulebox}[Regla cr\'itica de negocio]
No es posible registrar un traspaso SPEI si el proveedor tiene
\texttt{efos\_flag = True}, opini\'on SAT \texttt{vencida} o \texttt{no\_cargada},
o si la actividad econ\'omica no coincide. El sistema retorna un error de negocio
con mensaje espec\'ifico.
\end{rulebox}

\subsubsection{Entregables por Tipo de Proveedor}

\begin{table}[h!]
\centering
\begin{tabularx}{\textwidth}{>{\bfseries}lX}
\toprule
Tipo & Entregables adicionales \\
\midrule
Outsourcing & Lista de personal, turnos, asistencias firmadas, bit\'acoras
    operativas, evidencia fotogr\'afica, supervisi\'on documentada, relaci\'on de insumos \\[4pt]
Materiales & Fotos de materiales recibidos, acuse de recepci\'on firmado, registro de
    inventario, evidencia de entrega al destino final \\[4pt]
Servicios & Definici\'on de procesos, reportes con KPIs, evidencia de intervenci\'on
    (logs/tickets), resultados documentados, acta de entrega-recepci\'on \\
\bottomrule
\end{tabularx}
\caption{Entregables espec\'ificos por tipo de proveedor}
\end{table}

Todos los tipos comparten los \textbf{entregables base}: orden de compra firmada,
contrato vigente, responsable de comunicaci\'on y correos formales del per\'iodo.

\subsubsection{Traspasos SPEI}

El registro de un traspaso ejecuta \textbf{5 validaciones secuenciales}:

\begin{enumerate}
  \item Paso 2 (SAT) debe estar completo.
  \item La cuenta bancaria debe pertenecer al proveedor del folio.
  \item SAT vigente, EFOS negativo, actividad y RFC validados.
  \item Opini\'on SAT no vencida a la fecha.
  \item Monto $\leq$ presupuesto autorizado (o requiere aprobaci\'on de Direcci\'on).
\end{enumerate}

Si el monto supera el presupuesto y hay aprobaci\'on de Direcci\'on, el traspaso se
registra pero \textbf{genera autom\'aticamente una alerta} de tipo
\texttt{excede\_presupuesto} con severidad alta.

\subsubsection{CFDI}

Vincula comprobantes fiscales digitales a cada folio. Permite subir los archivos
XML y PDF del CFDI. Requiere que el Paso~4 (Traspaso) est\'e completo. Registra
UUID, monto, fecha de emisi\'on y estatus (vigente/cancelado). Las cancelaciones
deben incluir motivo documentado.

\subsubsection{Conciliaci\'on Bancaria}

Dos modalidades:

\begin{enumerate}
  \item \textbf{Carga de PDF}: El sistema analiza el estado de cuenta bancario con
    \texttt{pypdf} y extrae movimientos mediante expresiones regulares. Crea alertas
    autom\'aticas para cada anomal\'ia detectada.
  \item \textbf{Registro manual}: El operador captura movimientos individuales. Se
    aplican 4 reglas de alerta autom\'aticas:
    \begin{itemize}
      \item Dep\'osito sin folio $\rightarrow$ alerta alta \texttt{deposito\_sin\_folio}
      \item Monto $>$ presupuesto $\rightarrow$ alerta alta \texttt{excede\_presupuesto}
      \item Cuenta no autorizada $\rightarrow$ alerta alta \texttt{cuenta\_no\_autorizada}
      \item Sin CFDI vigente $\rightarrow$ alerta media \texttt{diferencia\_cfdi\_banco}
    \end{itemize}
\end{enumerate}

\subsubsection{Presupuestos}

Carga de presupuestos autorizados por per\'iodo (YYYY-MM). Dos modalidades:
carga masiva desde Excel con columnas \texttt{folio}, \texttt{presupuesto autorizado},
\texttt{periodo}; o alta manual individual por folio.

\subsubsection{Reportes Mensuales}

Generaci\'on autom\'atica de reporte ejecutivo con KPIs calculados:
\begin{itemize}
  \item Porcentaje de cierre de folios en el per\'iodo
  \item Alertas abiertas y alertas cr\'iticas
  \item N\'umero de proveedores con SAT vigente
  \item Ratio de cumplimiento SAT sobre total de proveedores
\end{itemize}

\subsubsection{Acuses de Conformidad}

Registro digital de la firma de conformidad de dos \'areas: UI y Contabilidad.
Incluye nombre del firmante, fecha y notas. Visible en el dashboard de Direcci\'on.

\newpage

% ═══════════════════════════════════════════════════════════════════════════════
\section{Responsabilidades Operativas por \'Area}
% ═══════════════════════════════════════════════════════════════════════════════

Esta secci\'on est\'a basada en la \textit{Carta de Instrucciones Operativas} emitida
por Direcci\'on General para UI y Contabilidad.

\subsection{Unidad de Inteligencia (UI)}

La UI es el \'area responsable de construir y mantener el expediente operativo de cada
folio SINGA. \textbf{Sin la evidencia generada por UI, no existe sustento para la
factura ni para el pago}.

\begin{table}[h!]
\centering
\begin{tabularx}{\textwidth}{Xlc}
\toprule
\textbf{Acci\'on} & \textbf{Plazo m\'aximo} & \textbf{KPI} \\
\midrule
Registrar folio SINGA en la plataforma & Mismo d\'ia de apertura & \% en tiempo \\
Subir orden de compra y contrato & 48 horas de iniciado & \% OC subidas \\
Subir evidencia operativa del per\'iodo & Antes del d\'ia 25 & \% expedientes completos \\
Subir asistencias y bit\'acoras & M\'ax.\ 5 d\'ias h\'abiles & \% asistencias \\
Subir evidencia fotogr\'afica & M\'ax.\ 5 d\'ias h\'abiles & \% fotos subidas \\
Comunicaciones relevantes del mes & Antes del cierre & \% comunicaciones \\
Marcar folio listo para revisi\'on & D\'ia 27 de cada mes & \% folios enviados \\
\bottomrule
\end{tabularx}
\caption{KPIs de la Unidad de Inteligencia}
\end{table}

\subsection{\'Area de Contabilidad}

\begin{table}[h!]
\centering
\begin{tabularx}{\textwidth}{Xlc}
\toprule
\textbf{Acci\'on} & \textbf{Plazo m\'aximo} & \textbf{KPI} \\
\midrule
Subir opini\'on SAT antes de autorizar pago & Previo al traspaso & \% opiniones previas \\
Subir Excel de presupuestos SINGA & D\'ia 5 de cada mes & \% meses con presupuesto \\
Correr conciliaci\'on bancaria & D\'ias 28--30 del mes & \% conciliaciones \\
Documentar y resolver alertas & 5 d\'ias h\'abiles & \% alertas en plazo \\
Vincular CFDIs a folios SINGA & M\'ax.\ 3 d\'ias de emitido & \% CFDIs vinculados \\
Entregar reporte mensual a Direcci\'on & D\'ia 5 del mes siguiente & \% reportes a tiempo \\
Verificar proveedores en lista EFOS & Antes del primer pago & \% verificados \\
\bottomrule
\end{tabularx}
\caption{KPIs del \'Area de Contabilidad}
\end{table}

\subsection{Conciliaciones Cr\'iticas Mensuales}

Contabilidad debe realizar y documentar \textbf{tres cruces de conciliaci\'on} antes
del cierre mensual:

\begin{table}[h!]
\centering
\begin{tabularx}{\textwidth}{lXl}
\toprule
\textbf{Conciliaci\'on} & \textbf{Qu\'e se cruza} & \textbf{Entregable} \\
\midrule
Bancos vs CFDIs vs Contabilidad &
  Cada dep\'osito/retiro debe tener CFDI y asiento contable. Sin diferencias. &
  Reporte conciliaci\'on IA \\[6pt]
IVA acreditable vs IVA trasladado &
  IVA de facturas recibidas vs IVA cobrado. &
  Reporte fiscal mensual \\[6pt]
N\'omina vs IMSS vs CFDIs &
  CFDIs de n\'omina coinciden con pagos SUA y n\'omina contabilizada. &
  Constancias SUA \\
\bottomrule
\end{tabularx}
\caption{Conciliaciones mensuales obligatorias}
\end{table}

\newpage

% ═══════════════════════════════════════════════════════════════════════════════
\section{Sistema de Alertas y Auditor\'ia}
% ═══════════════════════════════════════════════════════════════════════════════

\subsection{Centro de Alertas}

Las alertas se generan autom\'aticamente por el motor de conciliaci\'on y las
validaciones de traspasos. Cada alerta tiene severidad (\texttt{media}/\texttt{alta})
y estatus (\texttt{abierta}/\texttt{resuelta}).

\begin{rulebox}[Regla de resoluci\'on]
Toda alerta abierta debe resolverse en un plazo m\'aximo de \textbf{5 d\'ias h\'abiles}.
Las alertas de severidad alta que superen este plazo aparecen destacadas en el
dashboard de Direcci\'on. Un folio con alertas altas abiertas cambia autom\'aticamente
su estatus a \texttt{critico}.
\end{rulebox}

\subsection{Log de Auditor\'ia}

El modelo \texttt{AuditLog} registra \textbf{toda acci\'on importante} en el sistema.
La funci\'on \texttt{log\_action()} en \texttt{audit.py} se invoca desde todas las rutas
de escritura:

\begin{lstlisting}[language=Python, caption={Funci\'on de log de auditor\'ia}]
def log_action(*, user_id, action, entity, entity_id=None, details=None):
    """
    Crea un registro de auditoria en la base de datos.
    Solo usa argumentos nombrados (keyword-only).
    """
    entry = AuditLog(
        user_id=user_id,
        action=action,       # ej: 'folio_creado', 'traspaso_registrado'
        entity=entity,       # ej: 'Folio', 'Transfer', 'ProviderCompliance'
        entity_id=entity_id,
        details=details      # JSON o texto libre con contexto
    )
    db.session.add(entry)
    db.session.commit()
\end{lstlisting}

Eventos auditados: login/logout, creaci\'on de folios, actualizaci\'on de SAT,
registro de traspasos, vinculaci\'on de CFDI, carga de presupuestos, generaci\'on
de reportes.

\newpage

% ═══════════════════════════════════════════════════════════════════════════════
\section{C\'alculo de Fechas L\'imite y KPI de Entregables}
% ═══════════════════════════════════════════════════════════════════════════════

El m\'odulo \texttt{services.py} implementa la l\'ogica de c\'alculo de plazos y
el estado KPI de cada entregable.

\begin{lstlisting}[language=Python, caption={C\'alculo de fechas l\'imite por tipo de entregable}]
def due_date_for_code(code, owner_area):
    today = date.today()
    if code in ('oc', 'contrato'):
        return add_business_days(today, 2)
    elif code in ('asistencias', 'fotos', 'bitacoras', 'insumos',
                  'acuse_recepcion', 'inventario', 'evidencia_entrega',
                  'procesos', 'kpis', 'logs', 'resultados'):
        return add_business_days(today, 5)
    elif code == 'comunicaciones':
        return month_cutoff(25)     # dia 25 del mes en curso
    elif code in ('sat', 'efos'):
        return add_business_days(today, 1)
    elif code == 'cfdi':
        return add_business_days(today, 3)
    elif code == 'conciliacion':
        return month_cutoff(30)     # dia 30 del mes en curso
    return add_business_days(today, 5)

def kpi_status_for_item(item):
    if item.uploaded:
        return 'completo'
    if item.due_date and item.due_date < date.today():
        return 'vencido'
    return 'pendiente'
\end{lstlisting}

\newpage

% ═══════════════════════════════════════════════════════════════════════════════
\section{Datos de Demostraci\'on}
% ═══════════════════════════════════════════════════════════════════════════════

El archivo \texttt{app/seed.py} carga datos de demostraci\'on de forma
\textbf{idempotente} (no duplica registros si ya existen).

\subsection{Usuarios Demo}

\begin{table}[h!]
\centering
\begin{tabular}{lll}
\toprule
\textbf{Email} & \textbf{Rol} & \textbf{Contrase\~na} \\
\midrule
salo@batia.local        & Direcci\'on   & servicia2026 \\
mgonzalez@batia.local   & Tesorer\'ia   & servicia2026 \\
lhernandez@batia.local  & Tesorer\'ia   & servicia2026 \\
rfuentes@batia.local    & UI            & servicia2026 \\
cmorales@batia.local    & UI            & servicia2026 \\
pramirez@batia.local    & Contabilidad  & servicia2026 \\
\bottomrule
\end{tabular}
\caption{Usuarios de demostraci\'on}
\end{table}

\subsection{Proveedores y Folios Demo}

\begin{itemize}
  \item \textbf{Limpiadores SA} (outsourcing) --- con cuentas Banorte y BBVA
  \item \textbf{Insumos Alfa} (materiales) --- con cuenta BBVA
  \item \textbf{TechServ} (servicios) --- con cuenta Banorte
  \item \textbf{Mant.\ Delta} (outsourcing) --- con cuenta Banorte
\end{itemize}

Tres folios demo con estatus variables para prueba de todas las funcionalidades.

\newpage

% ═══════════════════════════════════════════════════════════════════════════════
\section{Despliegue y Operaci\'on}
% ═══════════════════════════════════════════════════════════════════════════════

\subsection{Entornos}

\begin{table}[h!]
\centering
\begin{tabularx}{\textwidth}{>{\bfseries}lXl}
\toprule
Entorno & Base de datos & Servidor \\
\midrule
Desarrollo & SQLite (\texttt{instance/servicia.db}) & \texttt{python run.py} \\
Producci\'on & PostgreSQL (Cloud SQL, GCP) & Gunicorn en Cloud Run \\
\bottomrule
\end{tabularx}
\end{table}

\subsection{Dockerfile}

\begin{lstlisting}[style=plain, caption={Dockerfile de producci\'on}]
FROM python:3.11-slim

ENV PYTHONDONTWRITEBYTECODE=1
ENV PYTHONUNBUFFERED=1
ENV PORT=8080

WORKDIR /app
COPY requirements.txt .
RUN pip install --no-cache-dir -r requirements.txt
COPY . .

CMD ["sh", "-c", \
  "gunicorn --bind 0.0.0.0:${PORT} \
            --workers 1 --threads 8 --timeout 120 wsgi:app"]
\end{lstlisting}

\subsection{Cloud Run --- Configuraci\'on}

\begin{table}[h!]
\centering
\begin{tabular}{ll}
\toprule
\textbf{Par\'ametro} & \textbf{Valor} \\
\midrule
Instancias m\'inimas & 0 (scale-to-zero) \\
Instancias m\'aximas & 3 \\
CPU                  & 1 vCPU \\
Memoria              & 512 MiB \\
Timeout              & 120 segundos \\
Cuenta de servicio   & cloud-run-mlh@\ldots.iam.gserviceaccount.com \\
Regi\'on             & us-central1 \\
\bottomrule
\end{tabular}
\caption{Configuraci\'on de Cloud Run}
\end{table}

\subsection{Variables de Entorno}

\begin{table}[h!]
\centering
\begin{tabular}{ll}
\toprule
\textbf{Variable} & \textbf{Prop\'osito} \\
\midrule
\texttt{SECRET\_KEY}          & Clave de sesiones Flask \\
\texttt{JWT\_SECRET\_KEY}     & Clave para firma de tokens JWT \\
\texttt{DATABASE\_URL}        & Cadena de conexi\'on (SQLite o PostgreSQL) \\
\texttt{COMPANY\_NAME}        & Nombre de la empresa para branding \\
\texttt{UPLOAD\_FOLDER}       & Ruta local para archivos subidos \\
\texttt{MAX\_CONTENT\_LENGTH} & Tama\~no m\'aximo de upload (16~MB) \\
\bottomrule
\end{tabular}
\caption{Variables de entorno}
\end{table}

\subsection{Instalaci\'on en Desarrollo}

\begin{lstlisting}[style=plain, caption={Pasos de instalaci\'on local}]
# 1. Instalar dependencias
pip install -r requirements.txt

# 2. Configurar entorno
cp .env.example .env
# Editar .env con los valores correspondientes

# 3. Iniciar aplicacion (auto-seed incluido)
python run.py

# 4. Acceder a la plataforma
# http://127.0.0.1:5000
\end{lstlisting}

\newpage

% ═══════════════════════════════════════════════════════════════════════════════
\section{Reglas Absolutas del Sistema}
% ═══════════════════════════════════════════════════════════════════════════════

Las siguientes reglas est\'an implementadas en c\'odigo como bloqueos duros y
no tienen excepciones operativas:

\begin{rulebox}[Reglas absolutas --- Aplicables a todas las \'areas]
\begin{enumerate}
  \item No se registra ning\'un folio sin el n\'umero de SINGA correspondiente.
  \item No se puede traspasar dinero a una cuenta que no est\'e autorizada por
    Direcci\'on en la plataforma.
  \item No se emite ni se paga ninguna factura sin contrato firmado previo.
  \item No se deduce ning\'un gasto sin materializaci\'on completa del servicio.
  \item No se paga a ning\'un proveedor sin Opini\'on de Cumplimiento SAT
    vigente y cargada.
  \item Todo debe estar conciliado: \textbf{Banco $\leftrightarrow$ CFDI
    $\leftrightarrow$ Contabilidad $\leftrightarrow$ Folio SINGA}.
  \item Cualquier diferencia entre el dep\'osito y el presupuesto SINGA requiere
    autorizaci\'on escrita de Direcci\'on antes de proceder.
  \item Las alertas generadas por el sistema deben resolverse en m\'aximo
    \textbf{5 d\'ias h\'abiles}.
\end{enumerate}
\end{rulebox}

\bigskip

\begin{alertbox}[Riesgo fiscal cr\'itico]
Pagar a un proveedor que aparezca en la lista EFOS (Empresas que Facturan
Operaciones Simuladas) \textbf{invalida autom\'aticamente la deducci\'on fiscal}
y genera responsabilidad solidaria para la empresa. Servicia bloquea este
escenario en c\'odigo: un proveedor con \texttt{efos\_flag = True} no puede
recibir traspasos.
\end{alertbox}

\newpage

% ═══════════════════════════════════════════════════════════════════════════════
\section{Dise\~no de Interfaz}
% ═══════════════════════════════════════════════════════════════════════════════

\subsection{Paleta de Colores}

\begin{table}[h!]
\centering
\begin{tabular}{>{\ttfamily}lll}
\toprule
\textbf{Variable} & \textbf{Uso} & \textbf{Hex} \\
\midrule
ink-900   & Fondo oscuro, textos primarios    & \#0b1736 \\
sky-500   & Acento principal, links, headers  & \#1c78c0 \\
teal-500  & Acento secundario, subt\'itulos   & \#0f8d8f \\
lime-500  & Estados ok, \'exito               & \#5f9d2f \\
amber-500 & Advertencias, estado pendiente    & \#c97a1a \\
rose-500  & Errores, estados cr\'iticos       & \#c94a56 \\
paper-100 & Fondo de tarjetas y panels        & \#f8fafc \\
\bottomrule
\end{tabular}
\caption{Sistema de colores de la interfaz}
\end{table}

\subsection{Componentes Principales}

\begin{itemize}
  \item \textbf{Dashboard de rol}: Hero personalizado con KPIs focalizados y
    accesos r\'apidos.
  \item \textbf{Master-Detail de folios}: Lista filtrable con panel de detalle
    contextual.
  \item \textbf{Workflow visual}: Indicadores de los 8 pasos con estado
    completado/pendiente.
  \item \textbf{Alertas flash}: Auto-cierre en 5.2 segundos con animaci\'on.
  \item \textbf{Efectos visuales}: Orbs flotantes animados en el fondo
    (CSS keyframes).
  \item \textbf{Reveal al scroll}: Intersection Observer para animar elementos
    al entrar en viewport.
\end{itemize}

\newpage

% ═══════════════════════════════════════════════════════════════════════════════
\section{Acuse de Recibo}
% ═══════════════════════════════════════════════════════════════════════════════

\vspace{1cm}

\begin{table}[h!]
\centering
\begin{tabularx}{\textwidth}{XX}
\toprule
\multicolumn{1}{c}{\textbf{Unidad de Inteligencia (UI)}} &
\multicolumn{1}{c}{\textbf{\'Area de Contabilidad}} \\
\midrule
& \\[1.8cm]
\underline{\hspace{6cm}} & \underline{\hspace{6cm}} \\
Nombre y firma del responsable & Nombre y firma del responsable \\[0.5cm]
Fecha: \underline{\hspace{3cm}} & Fecha: \underline{\hspace{3cm}} \\[0.5cm]
\bottomrule
\end{tabularx}
\end{table}

\vspace{2cm}

\begin{center}
{\small\color{gray}
Este documento es de uso interno y confidencial.\\
Servicia $\cdot$ M\'odulo de Aseguramiento Contable $\cdot$ Batia\\[0.2cm]
Generado a partir del an\'alisis del c\'odigo fuente\\
y la Carta de Instrucciones Operativas oficial.\\[0.2cm]
Versi\'on 2.0 $\cdot$ Abril 2026
}
\end{center}

\end{document}
